% English translation of chapter1.tex, lines 926--1140.
% Source: Methods of Algebra, Volume 2, commit
% 9a5803ff77dd3257484cb177f851a73770a59dd3 (CC BY 4.0).
% stable-unit-id: o014.aljabr2.chapter1.kan-extensions

% segment-id: o014.li2.chapter1.kan-extensions.h001
\section{Kan Extensions}\label{sec:Kan-extension}
\hypertarget{o014.aljabr2.chapter1.kan-extensions}{ }

% segment-id: o014.li2.chapter1.kan-extensions.p001
This section begins with the problem of extending a functor. Consider
categories $\mathcal{C}$, $\mathcal{D}$, $\mathcal{E}$ and functors $K$ and
$F$ as in the following diagram:
% segment-id: o014.li2.chapter1.kan-extensions.g001
\[\begin{tikzcd}
	\mathcal{C} \arrow[d, "K"'] \arrow[rd, bend left, "F"] & \\
	\mathcal{D} \arrow[dashed, r, "{\exists ? \; L}"'] & \mathcal{E}
\end{tikzcd}\]
In the basic spirit of category theory, we would like to find a functor $L$
indicated by the dashed arrow, together with a natural isomorphism
$F \simeq LK$. This problem generally has no solution: for example, there may
be $f \in \Mor(\mathcal{C})$ such that $Kf$ is an isomorphism while $Ff$ is
not. Failing that, we may at least ask whether there is a best approximation.
The answer is characterized by a universal property, with left and right
versions.

% segment-id: o014.li2.chapter1.kan-extensions.d001
\begin{definition}[D.\ Kan]\label{def:Kan-extension}
	\index{Kan extension}
	\index[sym1]{LanKF@$\Lan_K F$}
	\index[sym1]{RanKF@$\Ran_K F$}
	Consider categories $\mathcal{C}$, $\mathcal{D}$, $\mathcal{E}$ and
	functors $K: \mathcal{C} \to \mathcal{D}$ and
	$F: \mathcal{C} \to \mathcal{E}$.
% segment-id: o014.li2.chapter1.kan-extensions.l001
	\begin{itemize}
% segment-id: o014.li2.chapter1.kan-extensions.i001
		\item A \emph{left Kan extension} of the functor $F$ along $K$ is the
		following data $(\Lan_K F, \eta)$:
% segment-id: o014.li2.chapter1.kan-extensions.l002
		\begin{compactitem}
% segment-id: o014.li2.chapter1.kan-extensions.i002
			\item $\Lan_K F: \mathcal{D} \to \mathcal{E}$ is a functor;
% segment-id: o014.li2.chapter1.kan-extensions.i003
			\item $\eta: F \to (\Lan_K F) K$ is a natural transformation.
		\end{compactitem}
		These data must satisfy the following universal property: for every
		$L: \mathcal{D} \to \mathcal{E}$ and $\xi: F \to LK$, there is a unique
		natural transformation $\chi: \Lan_K F \to L$ such that
		$\xi = (\chi K) \eta$ (using vertical and horizontal composition of
		natural transformations), or, in $2$-cell diagrams:
% segment-id: o014.li2.chapter1.kan-extensions.g002
		\[\begin{tikzcd}[column sep=large]
			\mathcal{C} \arrow[d, "K"'] \arrow[rd, bend left, "F", ""' {name=F}] & \\
			\mathcal{D} \arrow[r, "L"'] \arrow[Rightarrow, from=F, "\xi" description] & \mathcal{E}
% segment-id: o014.li2.chapter1.kan-extensions.g003
		\end{tikzcd} = \begin{tikzcd}[column sep=large]
			\mathcal{C} \arrow[d, "K"'] \arrow[rd, bend left, "F", ""' {name=F}] & \\
			\mathcal{D} \arrow[r, "{\Lan_K F}"' name=A] \arrow[Rightarrow, from=F, "\eta" description] \arrow[r, "L"{below, name=B}, rounded corners, to path={-- ([yshift=-2em]\tikztostart.south) -- ([yshift=-2em] \tikztotarget.south) [midway]\tikztonodes -- (\tikztotarget) }] & \mathcal{E} \arrow[Rightarrow, from=A, to=B, "\chi"]
		\end{tikzcd}. \]

% segment-id: o014.li2.chapter1.kan-extensions.i004
		\item A \emph{right Kan extension} of the functor $F$ along $K$ is the
		following data $(\Ran_K F, \varepsilon)$:
% segment-id: o014.li2.chapter1.kan-extensions.l003
		\begin{compactitem}
% segment-id: o014.li2.chapter1.kan-extensions.i005
			\item $\Ran_K F: \mathcal{D} \to \mathcal{E}$ is a functor;
% segment-id: o014.li2.chapter1.kan-extensions.i006
			\item $\varepsilon: (\Ran_K  F)K \to F$ is a natural transformation.
		\end{compactitem}
		These data must satisfy the following universal property: for every
		$R: \mathcal{D} \to \mathcal{E}$ and $\delta: RK \to F$, there is a unique
		natural transformation $\theta: R \to \Ran_K F$ such that
		$\delta = \varepsilon (\theta K)$, or, in $2$-cell diagrams:
% segment-id: o014.li2.chapter1.kan-extensions.g004
		\[\begin{tikzcd}[column sep=large]
			\mathcal{C} \arrow[d, "K"'] \arrow[rd, bend left, "F", ""' {name=F}] & \\
			\mathcal{D} \arrow[r, "R"'] \arrow[Rightarrow, to=F, "\delta" description] & \mathcal{E}
% segment-id: o014.li2.chapter1.kan-extensions.g005
		\end{tikzcd} = \begin{tikzcd}[column sep=large]
			\mathcal{C} \arrow[d, "K"'] \arrow[rd, bend left, "F", ""' {name=F}] & \\
			\mathcal{D} \arrow[r, "{\Ran_K F}"' name=A] \arrow[Rightarrow, to=F, "\varepsilon" description] \arrow[r, "R"{below, name=B}, rounded corners, to path={-- ([yshift=-2em]\tikztostart.south) -- ([yshift=-2em] \tikztotarget.south) [midway]\tikztonodes -- (\tikztotarget) }] & \mathcal{E} \arrow[Rightarrow, from=B, to=A, "\theta"]
		\end{tikzcd}. \]
	\end{itemize}
\end{definition}


% segment-id: o014.li2.chapter1.kan-extensions.p002
If $\mathcal{C}$, $\mathcal{D}$, and $\mathcal{E}$ in the definition are
replaced respectively by $\mathcal{C}^{\opp}$, $\mathcal{D}^{\opp}$, and
$\mathcal{E}^{\opp}$, the pattern of directions of the functors remains the
same, while the natural transformations reverse direction. Thus left and
right Kan extensions are dual notions.

% segment-id: o014.li2.chapter1.kan-extensions.p003
Using the functor categories $\mathcal{E}^{\mathcal{C}}$ and
$\mathcal{E}^{\mathcal{D}}$, the universal properties of left and right Kan
extensions state, respectively, the bijections
% segment-id: o014.li2.chapter1.kan-extensions.q001
\begin{equation}\label{eqn:Kan-univ}\begin{aligned}
	\Hom_{\mathcal{E}^{\mathcal{D}}}\left( \Lan_K F, L \right) & \xrightarrow{1:1} \Hom_{\mathcal{E}^{\mathcal{C}}}\left( F, LK \right) \\
	\chi & \longmapsto (\chi K)\eta , \\
	\Hom_{\mathcal{E}^{\mathcal{D}}}\left( R, \Ran_K F \right) & \xrightarrow{1:1} \Hom_{\mathcal{E}^{\mathcal{C}}}\left( RK, F \right) \\
	\theta & \longmapsto \varepsilon (\theta K);
\end{aligned}\end{equation}
their inverses are the maps $(L, \xi) \mapsto \chi$ and
$(R, \delta) \mapsto \theta$, respectively, from
Definition~\ref{def:Kan-extension}.

% segment-id: o014.li2.chapter1.kan-extensions.pr001
\begin{proposition}\label{prop:Kan-uniqueness}
	Given categories $\mathcal{C}$, $\mathcal{D}$, $\mathcal{E}$ and functors
	$K: \mathcal{C} \to \mathcal{D}$ and $F: \mathcal{C} \to \mathcal{E}$,
	consider the pullback (that is, precomposition) functor
	$K^*: \mathcal{E}^{\mathcal{D}} \to \mathcal{E}^{\mathcal{C}}$ between
	the functor categories.
% segment-id: o014.li2.chapter1.kan-extensions.l004
	\begin{enumerate}[(i)]
% segment-id: o014.li2.chapter1.kan-extensions.i007
		\item If it exists, a left Kan extension $(\Lan_K F, \eta)$ is unique up
		to a unique isomorphism in $\mathcal{E}^{\mathcal{D}}$; the same holds for
		a right Kan extension.
% segment-id: o014.li2.chapter1.kan-extensions.i008
		\item If the left Kan extension (or right Kan extension) along $K$ exists
		for every $F$, these extensions assemble into a left adjoint to $K^*$,
		$\Lan_K: \mathcal{E}^{\mathcal{C}} \to \mathcal{E}^{\mathcal{D}}$
		(or a right adjoint
		$\Ran_K: \mathcal{E}^{\mathcal{C}} \to \mathcal{E}^{\mathcal{D}}$);
		the family of natural transformations assembled from $\eta$ (or
		$\varepsilon$) in Definition~\ref{def:Kan-extension} is exactly the unit
		(or counit) of this adjunction.
	\end{enumerate}
\end{proposition}
% segment-id: o014.li2.chapter1.kan-extensions.pf001
\begin{proof}
	For (i), simply apply the universal property in
	Definition~\ref{def:Kan-extension} in the usual way.

	For (ii), observe that $LK = K^*(L)$ and $RK = K^*(R)$; the desired
	adjunction is precisely the content of \eqref{eqn:Kan-univ}. The assertion
	about the unit or counit follows similarly; see
	\cite[Proposition 2.6.5]{Li1}.
\end{proof}

% segment-id: o014.li2.chapter1.kan-extensions.p004
By convention, the data $\eta$ or $\varepsilon$ in a Kan extension are often
omitted from the notation. If the left Kan extension (or right Kan extension)
along $K$ exists for every $F$, the functor $\Lan_K$ (or $\Ran_K$) obtained
in (ii) also has the corresponding uniqueness property, by uniqueness of
adjoints \cite[Proposition 2.6.10]{Li1}.

% segment-id: o014.li2.chapter1.kan-extensions.r001
\begin{remark}
	Since the definition of a Kan extension involves only composition of
	$2$-cells, it extends to general $2$-categories; the present situation is
	the special case of the $2$-category $\cate{Cat}$. See
	\cite[\S 3.5]{Li1} for details.
\end{remark}

% segment-id: o014.li2.chapter1.kan-extensions.p005
For every category $\mathcal{C}$, denote the unique functor
$\mathcal{C} \to \mathbf{1}$ by $!$; specifying a functor
$\mathbf{1} \to \mathcal{C}$ is equivalent to specifying an object of
$\mathcal{C}$. For every category $\mathcal{E}$ and object $X$ in it, write
$\Delta(X): \mathcal{C} \to \mathcal{E}$ for the composite
$\mathcal{C} \xrightarrow{!} \mathbf{1} \xrightarrow{\text{constant}\; X}
\mathcal{E}$; this is the constant functor that sends every object to $X$ and
every morphism to $\identity_X$.

% segment-id: o014.li2.chapter1.kan-extensions.e001
\begin{example}[Limits as Kan extensions]\label{eg:limit-as-Kan}
	Let $F: \mathcal{C} \to \mathcal{E}$ be a functor. Consider the diagram
% segment-id: o014.li2.chapter1.kan-extensions.g006
	\[\begin{tikzcd}
		\mathcal{C} \arrow[d, "!"'] \arrow[rd, bend left, "F"] & \\
		\mathbf{1} & \mathcal{E}
	\end{tikzcd}\]
	Then $\Lan_! F$ (or $\Ran_! F$) exists if and only if $\varinjlim F$ (or
	$\varprojlim F$) exists in $\mathcal{E}$, and in that case the Kan
	extension is precisely the limit in question.

	For example, consider $(\Lan_! F, \eta)$. The functor
	$\Lan_! F: \mathbf{1} \to \mathcal{E}$ may be viewed as an object of
	$\mathcal{E}$, while $\eta$ is a natural transformation
	$F \to \Delta(\Lan_! F)$ in $\mathcal{E}^{\mathcal{C}}$. Its universal
	property is equivalent to the assertion that the data $(\Lan_! F, \eta)$
	give an initial object of the category $(F/\Delta)$; this is precisely the
	description of $\varinjlim F$ in \S\ref{sec:limit-functor}.
\end{example}

% segment-id: o014.li2.chapter1.kan-extensions.e002
\begin{example}[Adjunctions as Kan extensions]
	Consider a pair of functors:
% segment-id: o014.li2.chapter1.kan-extensions.g007
	\[\begin{tikzcd}
		F: \mathcal{C} \arrow[r, shift left] & \mathcal{D}: G \arrow[l, shift left]
	\end{tikzcd}\]
	If $(F, G)$ is equipped with the structure of an adjunction with unit
	$\eta: \identity_{\mathcal{C}} \to GF$ and counit
	$\varepsilon: FG \to \identity_{\mathcal{D}}$, then $(G, \eta)$ gives
	$\Lan_F (\identity_{\mathcal{C}})$ and $(F, \varepsilon)$ gives
	$\Ran_G (\identity_{\mathcal{D}})$.

	To explain this, first we show that the following pullback functors
% segment-id: o014.li2.chapter1.kan-extensions.g008
	\[\begin{tikzcd}
		G^*: \mathcal{C}^{\mathcal{C}} \arrow[r, shift left] &
		\mathcal{C}^{\mathcal{D}}: F^* \arrow[l, shift left]
	\end{tikzcd}\]
	themselves form a natural adjoint pair. Observe that $\eta$
	induces
	$\identity_{\mathcal{C}}^* = \identity_{\mathcal{C}^{\mathcal{C}}}
	\to F^* G^* = (GF)^*$, and similarly $\varepsilon$ induces
	$G^* F^* \to \identity_{\mathcal{C}^{\mathcal{D}}}$; we continue to denote
	both by $\eta$ and $\varepsilon$. We must check the triangle identities
	characterizing the adjunction $(G^*, F^*, \eta, \varepsilon)$. In brief,
	these identities follow formally by applying pullback $(-)^*$ in functor
	categories to the triangle identities for $(F, G, \eta, \varepsilon)$;
	the details are left to the reader.

	This adjunction gives a canonical bijection
	$\Hom_{\mathcal{C}^{\mathcal{D}}}(\underbracket{G^* \identity_{\mathcal{C}}}_{= G}, L) \rightiso \Hom_{\mathcal{C}^{\mathcal{C}}}(\identity_{\mathcal{C}} , \underbracket{F^* L}_{=LF})$,
	for arbitrary $L: \mathcal{D} \to \mathcal{C}$; when $L = G$,
	$\identity_G$ maps to $\eta$. Writing the universal property of
	$\Lan_F (\identity_{\mathcal{C}})$ as the bijection in
	\eqref{eqn:Kan-univ}, one sees at once that $(G, \eta)$ gives
	$\Lan_F (\identity_{\mathcal{C}})$. The argument for
	$\Ran_G (\identity_{\mathcal{D}})$ is entirely similar.\footnote{Translator's
	note: the source writes $\Ran_F(\identity_{\mathcal{C}})$ in the final
	sentence, contrary to the preceding assertion concerning
	$(F,\varepsilon)$. The target corrects this to
	$\Ran_G(\identity_{\mathcal{D}})$ (O014-C014).}
\end{example}

% segment-id: o014.li2.chapter1.kan-extensions.p006
In practice, some situations require a further strengthening of the notion of
Kan extension\footnote{One common strengthening is called a pointwise Kan
extension; it is not discussed here.}, especially in the later treatment of
derived functors.

% segment-id: o014.li2.chapter1.kan-extensions.d002
\begin{definition}\label{def:absolute-Kan}
	Consider categories $\mathcal{C}$, $\mathcal{D}$, $\mathcal{E}$,
	$\mathcal{F}$ and functors $K: \mathcal{C} \to \mathcal{D}$ and
	$F: \mathcal{C} \to \mathcal{E}$.\footnote{Translator's note: the source
	does not include the category $\mathcal{F}$ in the initial list, even though
	it is used as the codomain of $M$. The target adds it (O014-C015).}
% segment-id: o014.li2.chapter1.kan-extensions.l005
	\begin{itemize}
% segment-id: o014.li2.chapter1.kan-extensions.i009
		\item Suppose the right Kan extension $(\Ran_K F, \varepsilon)$ exists.
		A functor $M: \mathcal{E} \to \mathcal{F}$ is said to
		\emph{preserve} $\Ran_K F$ if the composite $2$-cell
% segment-id: o014.li2.chapter1.kan-extensions.g009
		\[\begin{tikzcd}[column sep=large]
			\mathcal{C} \arrow[d, "K"'] \arrow[rd, bend left, "F", ""' {name=F}] & & \\
			\mathcal{D} \arrow[r, "{\Ran_K F}"'] \arrow[Rightarrow, to=F, "\varepsilon" description] & \mathcal{E} \arrow[r, "M"'] & \mathcal{F}
% segment-id: o014.li2.chapter1.kan-extensions.g010
		\end{tikzcd} =: \begin{tikzcd}
			\mathcal{C} \arrow[rr, "MF", ""{name=MF}] \arrow[rd, "K"'] & & \mathcal{F}  \\
			& \mathcal{D} \arrow[ru, "{M \Ran_K F}"'] \arrow[Rightarrow, to=MF, "M\varepsilon" description] &
		\end{tikzcd}\]
		gives a right Kan extension of $MF$ along $K$.
% segment-id: o014.li2.chapter1.kan-extensions.i010
		\item Suppose the left Kan extension $(\Lan_K F, \eta)$ exists. A functor
		$M: \mathcal{E} \to \mathcal{F}$ is said to
		\emph{preserve} $\Lan_K F$ if the composite $2$-cell
% segment-id: o014.li2.chapter1.kan-extensions.g011
		\[\begin{tikzcd}[column sep=large]
			\mathcal{C} \arrow[d, "K"'] \arrow[rd, bend left, "F", ""' {name=F}] & & \\
			\mathcal{D} \arrow[r, "{\Lan_K F}"'] \arrow[Rightarrow, from=F, "\eta" description] & \mathcal{E} \arrow[r, "M"'] & \mathcal{F}
% segment-id: o014.li2.chapter1.kan-extensions.g012
		\end{tikzcd} =: \begin{tikzcd}
			\mathcal{C} \arrow[rr, "MF", ""{name=MF}] \arrow[rd, "K"'] & & \mathcal{F}  \\
			& \mathcal{D} \arrow[ru, "{M \Lan_K F}"'] \arrow[Rightarrow, from=MF, "M\eta" description] &
		\end{tikzcd}\]
		gives a left Kan extension of $MF$ along $K$.
% segment-id: o014.li2.chapter1.kan-extensions.i011
		\item If a left Kan extension $(\Lan_K F, \eta)$ (or a right Kan
		extension $(\Ran_K F, \varepsilon)$) is preserved by every functor whose
		domain is $\mathcal{E}$, the Kan extension is called \emph{absolute}.
		\index{Kan extension!absolute}
	\end{itemize}
\end{definition}

% segment-id: o014.li2.chapter1.kan-extensions.p007
The next step is to explain the relationship between adjunctions and absolute
Kan extensions. The result is extremely useful in the later study of derived
functors; studying its proof also helps one become accustomed to Kan
extensions and the manipulation of $2$-cells. Consider a pair of adjoint
functors
% segment-id: o014.li2.chapter1.kan-extensions.g013
$\begin{tikzcd}
	F: \mathcal{C} \arrow[r, shift left] & \mathcal{C}' \arrow[l, shift left] : G
\end{tikzcd}$,
specified by a unit $\eta: \identity_{\mathcal{C}} \to GF$ and a counit
$\varepsilon: FG \to \identity_{\mathcal{C}'}$. Also given are functors
% segment-id: o014.li2.chapter1.kan-extensions.q002
\[ K: \mathcal{C} \to \mathcal{D}, \quad K': \mathcal{C}' \to \mathcal{D}'. \]
We shall now use $2$-cell diagrams systematically to represent functors,
natural transformations between them, and their composites, making the
argument more concise. In manipulating $2$-cells, we interchange vertical and
horizontal composition without further comment; this is valid, as explained
in \cite[Lemma 2.2.7]{Li1}.

% segment-id: o014.li2.chapter1.kan-extensions.th001
\begin{theorem}[{G.\ Maltsiniotis \cite{Ma07}}]\label{prop:Quillen-adjunction-gen}
	Suppose $K'F$ has an absolute right Kan extension along $K$, denoted
	$\mathrm{L}F$, while $KG$ has an absolute left Kan extension along $K'$,
	denoted $\mathrm{R}G$. The associated data are displayed in the following
	diagrams:
% segment-id: o014.li2.chapter1.kan-extensions.g014
	\[\begin{tikzcd}
		\mathcal{C} \arrow[r, "F"] \arrow[d, "K"'] & \mathcal{C}' \arrow[d, "{K'}"] \\
		\mathcal{D} \arrow[r, "{\mathrm{L}F}"'] \arrow[Rightarrow, ru, "\alpha" description] & \mathcal{D}'
% segment-id: o014.li2.chapter1.kan-extensions.g015
	\end{tikzcd} \quad \begin{tikzcd}
		\mathcal{C}' \arrow[r, "G"] \arrow[d, "{K'}"'] & \mathcal{C} \arrow[d, "K"] \arrow[Rightarrow, ld, "\beta" description] \\
		\mathcal{D}' \arrow[r, "{\mathrm{R}G}"'] & \mathcal{D} .
	\end{tikzcd} \]
	In this situation, there are unique transformations
	$\underline{\eta}: \identity_{\mathcal{D}} \to \mathrm{R}G \circ \mathrm{L}F$
	and
	$\underline{\varepsilon}: \mathrm{L}F \circ \mathrm{R}G \to
	\identity_{\mathcal{D}'}$ such that the following equations of composite
	$2$-cells hold:
% segment-id: o014.li2.chapter1.kan-extensions.q003
	\begin{equation*}\begin{gathered}
% segment-id: o014.li2.chapter1.kan-extensions.g016
		\begin{tikzcd}[every cell/.append style = {font = \small}]
			\mathcal{C} \arrow[r] \arrow[rr, bend left=75, "{\identity_{\mathcal{C}}}", ""' {name=I}] & \mathcal{C}' \arrow[r] \arrow[d] \arrow[Rightarrow, from=I, "\eta" description] & \mathcal{C} \arrow[d] \arrow[Rightarrow, ld, "\beta" description] \\
			& \mathcal{D}' \arrow[r, "{\mathrm{R}G}"'] & \mathcal{D}
% segment-id: o014.li2.chapter1.kan-extensions.g017
		\end{tikzcd} = \begin{tikzcd}[every cell/.append style = {font = \small}]
			\mathcal{C} \arrow[r] \arrow[d] & \mathcal{C}' \arrow[d] & \\
			\mathcal{D} \arrow[r, "{\mathrm{L}F}"'] \arrow[rr, bend right=75, "{\identity_{\mathcal{D}}}"', "" {name=I}] \arrow[Rightarrow, ru, "\alpha" description] & \mathcal{D}' \arrow[r, "{\mathrm{R}G}"'] \arrow[Rightarrow, from=I, "\underline{\eta}" description] & \mathcal{D}
		\end{tikzcd} \\
% segment-id: o014.li2.chapter1.kan-extensions.g018
		\begin{tikzcd}[every cell/.append style = {font = \small}]
			\mathcal{C}' \arrow[r] \arrow[bend left=75, rr, "{\identity_{\mathcal{C}'}}", ""' {name=I}] & \mathcal{C} \arrow[r] \arrow[d] \arrow[Rightarrow, to=I, "\varepsilon" description] & \mathcal{C}' \arrow[d] \\
			& \mathcal{D} \arrow[r, "{\mathrm{L}F}"'] \arrow[Rightarrow, ru, "\alpha" description] & \mathcal{D}'
% segment-id: o014.li2.chapter1.kan-extensions.g019
		\end{tikzcd} = \begin{tikzcd}
			\mathcal{C}' \arrow[r] \arrow[d] & \mathcal{C} \arrow[d] \arrow[Rightarrow, ld, "\beta" description] & \\
			\mathcal{D}' \arrow[r, "{\mathrm{R}G}"'] \arrow[rr, bend right=75, "{\identity_{\mathcal{D}'}}"', "" {name=I}] & \mathcal{D} \arrow[r, "{\mathrm{L}F}"'] \arrow[Rightarrow, to=I, "\underline{\varepsilon}" description] & \mathcal{D}'
		\end{tikzcd}
	\end{gathered}\end{equation*}
	Moreover, $\underline{\eta}$ and $\underline{\varepsilon}$ make
	$(\mathrm{L}F, \mathrm{R}G)$ an adjunction.
\end{theorem}
% segment-id: o014.li2.chapter1.kan-extensions.pf002
\begin{proof}
	%%%%%%%%%%%%%%%%%%%%%%%%%
	The first task is to prove the existence and uniqueness of
	$\underline{\eta}$ and $\underline{\varepsilon}$. Since $\mathrm{R}G$ is
	an absolute left Kan extension, the $2$-cell
% segment-id: o014.li2.chapter1.kan-extensions.g020
	\[\begin{tikzcd}
		\mathcal{C}' \arrow[r, "G"] \arrow[d, "{K'}"'] & \mathcal{C} \arrow[d, "K"] \arrow[Rightarrow, ld, "\beta" description] & \\
		\mathcal{D}' \arrow[r, "{\mathrm{R}G}"'] & \mathcal{D} \arrow[r, "{\mathrm{L}F}"'] & \mathcal{D}'
	\end{tikzcd} = (\mathrm{L}F)\beta : \mathrm{L}F \circ K \circ G \to \mathrm{L}F \circ \mathrm{R}G \circ K' \]
	makes $\mathrm{L}F \circ \mathrm{R}G$ a left Kan extension of
	$\mathrm{L}F \circ K \circ G$ along $K'$. On the other hand, we also have
% segment-id: o014.li2.chapter1.kan-extensions.g021
	\[\begin{tikzcd}[every cell/.append style = {font = \small}]
		\mathcal{C}' \arrow[r, "G"] \arrow[bend left=75, rr, "{\identity_{\mathcal{C}'}}", ""' {name=I}] & \mathcal{C} \arrow[r, "F"] \arrow[d, "K"'] \arrow[Rightarrow, to=I, "\varepsilon" description] & \mathcal{C}' \arrow[d, "{K'}"] \\
		& \mathcal{D} \arrow[r, "{\mathrm{L}F}"'] \arrow[Rightarrow, ru, "\alpha" description] & \mathcal{D}'
	\end{tikzcd} = \varepsilon(\alpha G) : \mathrm{L}F \circ K \circ G \to K' \circ \identity_{\mathcal{C}'} = K' . \]
	Thus the universal property of the left Kan extension determines a unique
	$\underline{\varepsilon}$ satisfying the equation of composite $2$-cells
	in the statement of the theorem.
	%%%%%%%%%%%%%%%%%%%%%%%%%

	The case of $\underline{\eta}$ is dual. It suffices to observe that
	$(\mathrm{R}G)\alpha:
	\mathrm{R}G \circ \mathrm{L}F \circ K \to
	\mathrm{R}G \circ K' \circ F$ also makes
	$\mathrm{R}G \circ \mathrm{L}F$ a right Kan extension, because
	$\mathrm{L}F$ is an absolute right Kan extension.

	It remains to verify the triangle identities for the unit
	$\underline{\eta}$ and counit $\underline{\varepsilon}$, namely the
	following equations of composite $2$-cells:
% segment-id: o014.li2.chapter1.kan-extensions.q004
	\begin{equation*}\begin{gathered}
% segment-id: o014.li2.chapter1.kan-extensions.g022
		\begin{tikzcd}[column sep=small, every cell/.append style = {font = \small}]
			\mathcal{D}' \arrow[r] \arrow[bend right=60, rr, "\identity"', "" {name=I}] & \mathcal{D} \arrow[r] \arrow[rr, bend left=60, "\identity", "" {name=J}] \arrow[Rightarrow, to=I] & \mathcal{D}' \arrow[r] \arrow[Rightarrow, from=J] & \mathcal{D}
% segment-id: o014.li2.chapter1.kan-extensions.g023
		\end{tikzcd} = \begin{tikzcd}[every cell/.append style = {font = \small}]
			\mathcal{D}' \arrow[r, bend left=60, "{\mathrm{R}G}", ""' {name=A}] \arrow[r, bend right=60, "{\mathrm{R}G}"', "" {name=B}] & \mathcal{D} \arrow[Rightarrow, from=A, to=B, "\identity" description]
		\end{tikzcd} , \quad
% segment-id: o014.li2.chapter1.kan-extensions.g024
		\begin{tikzcd}[column sep=small, every cell/.append style = {font = \small}]
			\mathcal{D} \arrow[r] \arrow[bend left=60, rr, "\identity", ""' {name=I}] & \mathcal{D}' \arrow[r] \arrow[Rightarrow, from=I] \arrow[bend right=60, rr, "\identity"', "" {name=J}] & \mathcal{D} \arrow[r] \arrow[Rightarrow, to=J] & \mathcal{D}'
% segment-id: o014.li2.chapter1.kan-extensions.g025
		\end{tikzcd} = \begin{tikzcd}[every cell/.append style = {font = \small}]
			\mathcal{D} \arrow[r, bend left=60, "{\mathrm{L}F}", ""' {name=A}] \arrow[r, bend right=60, "{\mathrm{L}F}"', "" {name=B}] & \mathcal{D}' \arrow[Rightarrow, from=A, to=B, "\identity" description]
		\end{tikzcd},
	\end{gathered}\end{equation*}
	where unambiguous arrow labels have been omitted. We prove the first
	equation. By the universal property of $(\mathrm{R}G, \beta)$ as a left Kan
	extension, it suffices to show that the $2$-cells obtained by composing both
	sides with $\beta$ are equal. For the left side of the first equation, begin
	with
% segment-id: o014.li2.chapter1.kan-extensions.g026
	\[\begin{tikzcd}[every cell/.append style = {font = \small}]
		\mathcal{C}' \arrow[r] \arrow[d] & \mathcal{C} \arrow[d] \arrow[Rightarrow, ld] & & \\
		\mathcal{D}' \arrow[r] \arrow[bend right=60, rr, "\identity"', "" {name=I}] & \mathcal{D} \arrow[r] \arrow[rr, bend left=60, "\identity", "" {name=J}] \arrow[Rightarrow, to=I, "\underline{\varepsilon}"] & \mathcal{D}' \arrow[r] \arrow[Rightarrow, from=J, "\underline{\eta}"] & \mathcal{D}
% segment-id: o014.li2.chapter1.kan-extensions.g027
	\end{tikzcd} = \begin{tikzcd}[every cell/.append style = {font = \small}]
		\mathcal{C}' \arrow[r] \arrow[rr, bend left=40, "\identity", ""' {name=T}] & \mathcal{C} \arrow[r] \arrow[d] \arrow[Rightarrow, to=T, "\varepsilon" inner sep=0.4em] & \mathcal{C}' \arrow[d] & \\
 		& \mathcal{D} \arrow[r] \arrow[rr, bend right=60, "\identity"', "" {name=I}] \arrow[Rightarrow, ru] & \mathcal{D}' \arrow[r] \arrow[Rightarrow, from=I, "\underline{\eta}"'] & \mathcal{D}
	\end{tikzcd}\]
	where we used the characterization of $\underline{\varepsilon}$. Next,
	using the characterization of $\underline{\eta}$ and the triangle
	identities satisfied by $\eta$ and $\varepsilon$, the right-hand side above
	becomes
% segment-id: o014.li2.chapter1.kan-extensions.g028
	\[\begin{tikzcd}[every cell/.append style = {font = \small}]
		\mathcal{C}' \arrow[r] \arrow[rr, bend right=75, "\identity"', "" {name=A}] & \mathcal{C} \arrow[r] \arrow[rr, bend left=75, "\identity", ""' {name=B}] \arrow[Rightarrow, to=A, "\varepsilon"] & \mathcal{C}' \arrow[r] \arrow[d] \arrow[Rightarrow, from=B, "\eta"] & \mathcal{C} \arrow[d] \arrow[Rightarrow, ld, "\beta" description] \\
		& & \mathcal{D}' \arrow[r] & \mathcal{D}
% segment-id: o014.li2.chapter1.kan-extensions.g029
	\end{tikzcd} = \begin{tikzcd}[column sep=large, every cell/.append style = {font = \small}]
		\mathcal{C}' \arrow[r] \arrow[bend left=75, r, ""' {name=T}] \arrow[r, "" {name=B}] \arrow[d] & \mathcal{C} \arrow[d] \arrow[Rightarrow, ld, "\beta" description] \\
		\mathcal{D}' \arrow[r] & \mathcal{D} \arrow[Rightarrow, from=T, to=B, "\identity"]
% segment-id: o014.li2.chapter1.kan-extensions.g030
	\end{tikzcd} = \begin{tikzcd}[column sep=large, every cell/.append style = {font = \small}]
		\mathcal{C}' \arrow[r] \arrow[d] & \mathcal{C} \arrow[d]  \arrow[Rightarrow, ld, "\beta" description] \\
		\mathcal{D}' \arrow[r, ""' {name=T}] \arrow[r, bend right=75, "" {name=B}] & \mathcal{D} \arrow[Rightarrow, from=T, to=B, "\identity"]
	\end{tikzcd}\]
	This is the first triangle identity. The proof of the second is dual.
\end{proof}
